%!TeX root = ../main.tex

\subsection{Análisis de estabilidad del sistema a través de la matriz dinámica}
Para este análisis se realizará un determinante y la
obtención de los autovalores con el fin de conocer los puntos de estabilidad
del sistema.

\subsubsection{Determinante de $\mathbf{A}$}

\begin{align*}
  det\left(\mathbf{A}\right) &=
  det\begin{bmatrix}
  -\frac{R}{L}          & 0        & -\frac{K_e}{L}   \\
  0                     & 0        &  1               \\
  \frac{K_t}{J}         & 0        & -\frac{B}{J}     \\
  \end{bmatrix} \\
\end{align*}
\begin{align*}
  %%% det(A)
  &det\left(\mathbf{A}\right) =\\
%  \left(-\tfrac{R}{L} \right)
%  \begin{vmatrix}
%    0              &  1           \\
%    0              & -\frac{B}{J} \\
%  \end{vmatrix}
%  -\left(0\right)
%  \begin{vmatrix}
%    0              &  1           \\
%    \frac{K_t}{J}  & -\frac{B}{J} \\
%  \end{vmatrix}
%  +\left(-\tfrac{K_e}{L}\right)
%  \begin{vmatrix}
%    0              & 0            \\
%    \frac{K_t}{J}  & 0            \\
%  \end{vmatrix}\\
  &\begin{bmatrix}
    -\tfrac{R}{L} 
    -0
    +-\tfrac{K_e}{L}
  \end{bmatrix}
  \begin{bmatrix}
    \begin{vmatrix}
      0              &  1           \\
      0              & -\frac{B}{J} \\
    \end{vmatrix}\\
    \begin{vmatrix}
      0              &  1           \\
      \frac{K_t}{J}  & -\frac{B}{J} \\
    \end{vmatrix}\\
    \begin{vmatrix}
      0              & 0            \\
      \frac{K_t}{J}  & 0            \\
    \end{vmatrix}
  \end{bmatrix}\\
  %%% det(A)
  %%% det(A)
  &det\left(\mathbf{A}\right) = 0
\end{align*}

Dicho determinante implica que el sistema tiene infinitos puntos
de estabilidad.

\subsubsection{Autovalores de $\mathbf{A}$}
Los autovalores son aquellos valores $\lambda$ que cumplen con
la ecuación \eqref{eq:DefinicionAutovalores}.
\begin{equation}
  det\left(\mathbf{A-\lambda I}\right)=0
  \label{eq:DefinicionAutovalores}
\end{equation}

Por lo tanto se necesita obtener la matriz \eqref{eq:Matriz A - lambda I}.

\begin{equation}
  \mathbf{A-\lambda I} =
  \begin{bmatrix}
  -\frac{R}{L}          & 0        & -\frac{K_e}{L}   \\
  0                     & 0        &  1               \\
  \frac{K_t}{J}         & 0        & -\frac{B}{J}     \\
  \end{bmatrix}
  \label{eq:Matriz A - lambda I}
\end{equation}

luego, se obtiene el determinante de \eqref{eq:Matriz A - lambda I}
que resulta ser un polinomio de tercer grado \eqref{eq:detA - lI}

\begin{align*}
  det\left(\mathbf{A-\lambda I}\right) &=
  det\begin{bmatrix}
  -\frac{R}{L}-\lambda  & 0        & -\frac{K_e}{L}         \\
  0                     & -\lambda &  1                     \\
  \frac{K_t}{J}         & 0        & -\frac{B}{J} - \lambda \\
  \end{bmatrix} \\
  %%% det(A)
  det\left(\mathbf{A-\lambda I}\right) &=
  \left(-\tfrac{R}{L} - \lambda\right)
  \begin{vmatrix}
    -\lambda       &  1                     \\
    0              & -\frac{B}{J} - \lambda \\
  \end{vmatrix}\\
  &-\left(0\right)
  \begin{vmatrix}
    0              &  1                     \\
    \frac{K_t}{J}  & -\frac{B}{J} - \lambda \\
  \end{vmatrix}\\
  &+\left(-\tfrac{K_e}{L}\right)
  \begin{vmatrix}
    0              & -\lambda               \\
    \frac{K_t}{J}  & 0                      \\
  \end{vmatrix}
\end{align*}

\begin{equation}
  \begin{aligned}
    %%% det(A)
    &det\left(\mathbf{A-\lambda I}\right) = \\
    &\left(-\tfrac{R}{L}-\lambda\right)\left(\tfrac{B}{J}\lambda+\lambda^2\right)
    -\left(\tfrac{K_e K_t}{JL}\lambda\right)
    \label{eq:detA  - lI}
  \end{aligned}
\end{equation}

Luego de obtener \eqref{eq:detA  - lI}, a fin de rescatar los autovalores se procede
a resolver la expresión \eqref{eq:DefinicionAutovalores} para la matriz \eqref{eq:Matriz A - lambda I}

\begin{equation*}
  \begin{aligned}
    det\left(\mathbf{A-\lambda I}\right) &=0  \\
    \left(-\tfrac{R}{L}-\lambda\right)\left(\tfrac{B}{J}\lambda+\lambda^2\right)
    -\left(\tfrac{K_e K_t}{JL}\lambda\right)
                                         &=0  \\
    \lambda\left[
    \left(-\tfrac{R}{L}-\lambda\right)\left(\tfrac{B}{J}+\lambda\right)
    -\tfrac{K_e K_t}{JL}
    \right]
                                         &=0  \\
    -\lambda\left[
    \left(\tfrac{R}{L}+\lambda\right)\left(\tfrac{B}{J}+\lambda\right)
    +\tfrac{K_e K_t}{JL}
    \right]
                                         &=0  \\
  \end{aligned}
\end{equation*}

La primera solución se puede encontrar de la forma mostrada en \eqref{eq:SolucionAutovalores 1}

Definiendo $p= \left(\tfrac{R}{L}+\lambda\right)\left(\tfrac{B}{J}+\lambda\right)
    +\tfrac{K_e K_t}{JL} $
\begin{equation}
  \begin{aligned}
    \frac1{-p}\left(-\lambda p\right)
     &=\frac1{-p}0\\
    \lambda_1 &=0\\
     \label{eq:SolucionAutovalores 1}
  \end{aligned}
\end{equation}

De manera análoga es posible encontrar los segundo y tercer autovalores
obteniendo la igualdad respecto a $p$ y resolviendo la ecuación
de segundo grado \eqref{eq:SolucionAutovalores 2} que se obtiene,
resultando en la expresión \eqref{eq:SolucionAutovalores 3}.

\begin{equation*}
  \begin{aligned}
    \frac1{-\lambda}\left(-\lambda p\right)
     &=\frac1{-\lambda}0\\
    p &=0\\
  \end{aligned}
\end{equation*}

\begin{equation}
  \begin{aligned}
    \left(\tfrac{R}{L}+\lambda\right)\left(\tfrac{B}{J}+\lambda\right)
    +\tfrac{K_e K_t}{JL} &=0
     \label{eq:SolucionAutovalores 2}
  \end{aligned}
\end{equation}

\begin{equation*}
      \lambda^2 + \left(\tfrac{B}{J} + \tfrac{R}{L}\right)\lambda
    +\tfrac{BR+K_e K_t}{JL} =0
\end{equation*}
\begin{equation}
  \begin{aligned}
    &\lambda_{2,3} = \\
    &\frac{
      - \left( \tfrac{B}{J} + \tfrac{R}{L} \right)
      \pm \sqrt{
        \left( \tfrac{B}{J} + \tfrac{R}{L} \right)^2
        - 4 \tfrac{BR+K_e K_t}{JL}
      }
    }
    {2}
  \end{aligned}
     \label{eq:SolucionAutovalores 3}
\end{equation}

\subsubsection{Análisis de autovalores con valores experimentales}

A partir de valores experimentales indicados en \cite{Ugurlu2025} para un motor de DC, se obtienen autovalores numéricos que ilustran la dinámica del sistema. Los parámetros utilizados como ejemplo son: resistencia $R = 2.1\ \Omega$, inductancia $L = 5\times10^{-3}\ \text{H}$, constante de par $K_t = 0.12\ \text{N·m/A}$, constante de fuerza contraelectromotriz $K_e = 0.12\ \text{V/(rad/s)}$, momento de inercia $J = 1\times10^{-4}\ \text{kg·m}^2$ y coeficiente de fricción viscosa $B = 1\times10^{-5}\ \text{N·m·s/rad}$.

Sustituyendo estos valores en la expresión general para $\lambda_{2,3}$ de \eqref{eq:SolucionAutovalores 3}, y definiendo  
\[
\delta = \left(\tfrac{B}{J} + \tfrac{R}{L}\right)^2 - 4 \tfrac{BR+K_e K_t}{JL},
\]  
se calcula:

\begin{equation*}
\begin{aligned}
\tfrac{B}{J} + \tfrac{R}{L} &= \tfrac{1\mathrm{e}{-5}}{1\mathrm{e}{-4}} + \tfrac{2.1}{5\mathrm{e}{-3}} \\
&= 1\mathrm{e}{-1} + 4.2\mathrm{e}{2} = 4.201\mathrm{e}{2} \\[2pt]
\left(\tfrac{B}{J} + \tfrac{R}{L}\right)^2 &= (4.201\mathrm{e}{2})^{2} = 1.76484\mathrm{e}{5} \\[2pt]
4 \tfrac{BR+K_e K_t}{JL} &= 4 \cdot \tfrac{(2.1\mathrm{e}{0})(1\mathrm{e}{-5}) + (1.2\mathrm{e}{-1})(1.2\mathrm{e}{-1})}{(1\mathrm{e}{-4})(5\mathrm{e}{-3})} \\
&= 4 \cdot \tfrac{2.1\mathrm{e}{-5} + 1.44\mathrm{e}{-2}}{5\mathrm{e}{-7}} \\
&= 4 \cdot \tfrac{1.4421\mathrm{e}{-2}}{5\mathrm{e}{-7}} \\
&= 4 \cdot 2.8842\mathrm{e}{4} = 1.15368\mathrm{e}{5} \\[2pt]
\delta &= 1.76484\mathrm{e}{5} - 1.15368\mathrm{e}{5}\\
       &= 6.1116\mathrm{e}{4} \\[2pt]
\sqrt{\delta} &\approx 2.4722\mathrm{e}{2}
\end{aligned}
\end{equation*}
Luego, los autovalores no nulos resultan:

\begin{equation*}
\begin{aligned}
\lambda_{2,3} &= \frac{ - \left( \tfrac{B}{J} + \tfrac{R}{L} \right) \pm \sqrt{\delta}}{2} \\[2pt]
\lambda_2 &\approx \frac{-4.201\mathrm{e}{2} + 2.4722\mathrm{e}{2}}{2} \approx  86.44 \\[2pt]
\lambda_3 &\approx \frac{-4.201\mathrm{e}{2} - 2.4722\mathrm{e}{2}}{2} \approx -333.66
\end{aligned}
\end{equation*}

De este modo, el conjunto completo de autovalores para este caso de ejemplo es:
\begin{equation}
    \lambda_1 = 0,\quad \lambda_2 \approx -86.44,\quad \lambda_3 \approx -333.66
    \label{eq:ResultadoAutovaloresExperimental}
\end{equation}

\subsubsection{Conclusión del análisis de estabilidad por medio de la matriz dinámica}

Debido a que la matriz dinámica tiene un autovalor $\lambda={0}$, se puede
concluir que el sistema es marginalmente estable.
